\documentclass[11pt,a4paper]{article}
\usepackage[T1]{fontenc}
\usepackage{lmodern,amsmath,amssymb,amsthm,booktabs,microtype,xcolor}
\usepackage[margin=24mm]{geometry}
\usepackage[bookmarks=false]{hyperref}
\definecolor{ink}{HTML}{244B68}
\hypersetup{colorlinks=true,linkcolor=ink,urlcolor=ink,pdftitle={BundleFuse: A Technical Publication Argument}}
\newtheorem{proposition}{Proposition}
\newcommand{\I}{\mathcal I}
\newcommand{\off}{\operatorname{off}}
\newcommand{\diag}{\operatorname{diag}}
\newcommand{\pos}[1]{\left[#1\right]_+}
\newcommand{\Z}{\mathcal Z_S}
\setlength{\parindent}{0pt}
\setlength{\parskip}{5pt}
\setlength{\emergencystretch}{2em}
\title{\textcolor{ink}{BundleFuse: A Technical Publication Argument}\\[5pt]\large Definitions, exact identities, and the evidence supporting them}
\author{}\date{5 November 2026}
\begin{document}
\maketitle
\vspace{-20pt}
\textbf{Central claim.} Historical item relations and evidence shared through item
attributes provide complementary estimates of bundle compatibility. BundleFuse
constructs and combines these estimates explicitly. Candidate-frequency backoff
and multi-item terms are additional modeling choices whose benefits must be
established empirically.

\textbf{Scope.} The equations below describe the binary-incidence, full-feature
research formulation used for the four supplied POG/Spotify splits. They explain
the common construction and its refinements; they are not an exhaustive specification
of every modality or of the retail/commerce variants. Numerical findings in
Section~\ref{evidence} use only the two completed POG mechanism studies available
when this note was revised. The music mechanism studies are pending.

\section{The prediction problem and the starting estimator}
Let $\I=\{1,\ldots,n\}$ be a fixed catalogue with $n>1$ items. Let
$S\subset\I$ be a nonempty observed set and let $c\in\I\setminus S$ be a candidate.
The objective is to assign a real ranking score $\theta(c\mid S)$ to each candidate.
A larger score indicates a preferred completion; the score is not defined as a
probability. Hidden completion targets are available for validation and evaluation,
but are not arguments of the scoring function.

Let $m$ be the number of training bundles and $B\in\{0,1\}^{m\times n}$ their incidence
matrix: $B_{ei}=1$ precisely when training bundle $e$ contains item $i$. Define
\[
 x_S(i)=\mathbf1[i\in S],\qquad f_i=\sum_{e=1}^mB_{ei},\qquad
 d_e=\sum_{i=1}^nB_{ei}.
\]
Here $\mathbf1[\cdot]$ is the indicator of a statement, $f_i$ counts training
appearances, and $d_e$ is bundle size. An elementary score is
\begin{equation}\label{raw}
 r_{\rm count}(c\mid S)=(B^\top Bx_S)_c
 =\sum_{s\in S}\sum_{e=1}^m B_{ec}B_{es}.
\end{equation}
For $c\notin S$, this counts recorded co-memberships with the observed items.
It motivates two modifications. First, observed counts can be dominated by frequent
items and large bundles. Second, a zero count does not distinguish incompatibility
from a lack of observations. These are reasons to consider modifications, not proofs
that any particular modification improves ranking.

\section{Construct historical estimates and extend their coverage}
\subsection{Weighted historical overlap}
Define diagonal matrices $D\in\mathbb R^{n\times n}$ and
$V\in\mathbb R^{m\times m}$ by
\[
 D_{ii}=(f_i+5)^{-1/2},\qquad V_{ee}=(d_e+1)^{-1/2}.
\]
The basic weighted overlap vector is
\begin{equation}\label{pair}
 p(S)=B^\top VBDx_S.
\end{equation}
For equal co-membership counts, common observed items and large training bundles
contribute less. $D$ acts on the observed-item side of this expression. Thus the
operator in~\eqref{pair} is not generally symmetric. The constants $5$, $1$ and the
exponents specify this version of the method; they are not derived optimal values.

For propagation, define $E_{ee}=(d_e+1)^{-1}$ and
\begin{equation}\label{propagation}
 L=\off(DB^\top EBD),
\end{equation}
where $\off(M)$ replaces a square matrix's diagonal by zero. The propagated score
is $Lp(S)$. This uses indirect historical relationships. $L$ is symmetric, but is
not in general row-stochastic. Equation~\eqref{pair} itself retains its diagonal.
Although that diagonal does not contribute directly to an eligible candidate's
score, it affects observed-item entries used in propagation and standardization.
It must not be silently removed from the definition of $p$.

Raw co-occurrence, positive pointwise mutual information, and collaborative singular
value decomposition provide other summaries of historical relations. They reuse
history and need not be statistically independent. Their inclusion is a choice of
candidate estimators; their necessity is an empirical question.

\subsection{Relations shared through categories and descriptions}
Let $g_i\in\{1,\ldots,q\}$ be item $i$'s category, and let
$A\in\{0,1\}^{n\times q}$ satisfy $A_{ia}=\mathbf1[g_i=a]$. Define
$T\in\mathbb R^{q\times q}$ so that $T_{ab}$ scores candidate category $a$ against
observed category $b$. A category-based score has the form
\begin{equation}\label{category}
 r_{\rm category}(S)=ATA^\top x_S.
\end{equation}
Pooling training observations across items can supply support when an exact pair in
\eqref{raw} has never occurred. Signed diagonal entries of $T$ can encode additive
penalties for repeated categories. The implemented category table and penalties are
specified by the domain;~\eqref{category} is their operator form, not a universal
formula for estimating $T$. Such a relation need not be a similarity or a
positive-semidefinite kernel.

Let $F\in\mathbb R^{n\times d}$ contain fixed item feature rows $F_i$. A descriptive
score used in the full-ranking path is
\begin{equation}\label{feature}
 r_F(S)=F\left(\frac{1}{|S|}\sum_{s\in S}F_s\right)^\top
       =\frac{1}{|S|}FF^\top x_S.
\end{equation}
The feature dimension $d$ and the row construction depend on the modality. Features
can distinguish candidates absent from the training bundles. Resemblance is only
one compatibility cue: distinct product categories may complement one another,
whereas artist or album repetition may be useful in music.

For a supplied feature row $X_i$, let $\ell_i$ be its centering-group label.
Fashion uses $\ell_i=g_i$; music uses labels obtained by clustering feature vectors,
not artist labels. For each nonempty group $a$, define
$\bar X_a=|\{i:\ell_i=a\}|^{-1}\sum_{i:\ell_i=a}X_i$.
A residual feature starts from $X_i-\bar X_{\ell_i}$ and then normalizes each row.
The residuals have zero group mean \emph{before} row normalization. The normalized
rows need not retain zero group mean, and this operation does not establish
independence from category or group identity.

To state the feature bridge, define row normalization
$\mathcal N(v)=v/(\lVert v\rVert_2+10^{-8})$, applied separately to each row of a
matrix. For this construction, let $F$ be the row-normalized concatenation of the
modality residuals in fashion, or the normalized audio residuals in music.
Its graph-adjusted form is
\begin{equation}
 \widetilde F=\mathcal N\bigl(F+\tfrac12\mathcal N(LF)\bigr).
\end{equation}
It mixes descriptions with descriptions propagated through historical relations.
Category relations and this bridge therefore use both attributes and history. The
two evidence families are functional groupings, not disjoint information sources.

\section{Define the combined score before interpreting it}
\subsection{Standardization and the additive core}
Let $\mathcal H$ be the available signal index set and let
$r_h(S)\in\mathbb R^n$ denote raw signal $h\in\mathcal H$. Define
\begin{align}
 a_h(S)&=\frac1n\sum_{i=1}^n r_h(i\mid S),\\
 b_h(S)&=\left[\frac{1}{n-1}\sum_{i=1}^n
          (r_h(i\mid S)-a_h(S))^2\right]^{1/2},\\
 z_h(c\mid S)&=\Z[r_h]_c
 =\frac{r_h(c\mid S)-a_h(S)}{b_h(S)+\varepsilon},\qquad \varepsilon=10^{-9}.
\end{align}
The divisor $n-1$ specifies the Torch full-ranking path considered here. Some
NumPy basket variants use divisor $n$ and are outside this numerical specification.
All catalogue items, including $S$, enter these statistics; items in $S$ are excluded
when recommendations are ranked. These operations depend on $S$ and fixed model
artifacts, not on the query's hidden targets. This establishes a dependency property
of inference, not an audit of upstream feature training or dataset provenance.

Let $w_h\geq0$ be a validation-selected weight. Without multi-item refinements,
let $z_h^0$ denote each standardized signal and define
\begin{equation}\label{core}
 Y_0(c\mid S)=\sum_{h\in\mathcal H}w_hz_h^0(c\mid S).
\end{equation}
Standardization makes the raw scales comparable; it does not make the signals
independent or equally reliable. Because its mean and denominator depend on $S$,
\eqref{core} need not equal a single fixed matrix times $x_S$.

\subsection{Precisely defined multi-item refinements}
We next define the optional corrections so that the eventual full score is explicit.
For each training bundle, let $k_e(S)=\sum_{s\in S}B_{es}$ and
$o_e(S)=V_{ee}\sum_{s\in S}B_{es}D_{ss}$. Define the incidence surplus
\begin{equation}
 \delta_c(S)=\sum_{e=1}^mB_{ec}o_e(S)\,\beta\pos{k_e(S)-1},
 \qquad \beta=1.5,\qquad \pos{x}=\max(x,0).
\end{equation}
Only training bundles containing at least two observed items contribute. Replace
$p$ by $p+\delta$ and $Lp$ by $Lp+L\delta$ in their respective raw channels, leaving
other channels unchanged. Write $z_h^1$ for the separately standardized modified
channels and $Y_1=\sum_h w_hz_h^1$.

For the second refinement, let $W=\off(B^\top VBD)$ and let $v_i$ be the selected
normalized descriptor row for item $i$: the title residual in POG and the audio
residual in Spotify. The category label $g_i$ used below is distinct from the music
centering-group label $\ell_i$. Define
candidate-to-seed support and seed similarity, respectively, by
\[
 q_{cs}=W_{cs}+0.15\pos{v_c\cdot v_s},\qquad
 G_{st}=\tfrac12\mathbf1[g_s=g_t]+\tfrac12\pos{v_s\cdot v_t}.
\]
Set $K_{ss}=0$. For distinct seeds $s,t$, use $K_{st}=\pos{1-G_{st}}$ for fashion
and $K_{st}=1+G_{st}$ for music. The raw Rao-based score and its standardized version are
\begin{equation}
 R(c\mid S)=\frac{\sum_{s,t\in S}q_{cs}K_{st}q_{ct}}
                    {\sum_{s\in S}q_{cs}+0.05},\qquad z_R=\Z[R].
\end{equation}
This evaluates support involving pairs of observed items under the chosen category/
descriptor relationship. Both $\delta$ and $R$ vanish when $|S|=1$.

For selected $\alpha\in[0,1]$ and $\gamma\geq0$, define the full combined estimate
\begin{equation}\label{Y}
 Y=(1-\alpha)Y_0+\alpha Y_1+\gamma z_R.
\end{equation}
The mixture is formed \emph{after} each raw channel is standardized. In general,
$\Z[p+\alpha\delta]\ne(1-\alpha)\Z[p]+\alpha\Z[p+\delta]$.
These corrections expand the construction; they are not assumed to improve every dataset.

\subsection{The restricted fallback and its relation to the full estimate}
Let $\mathcal Q\subset\mathcal H$ be the available fallback-eligible channels:
category, soft-category, text-style, image-style and bridge for fashion; category,
group homophily, audio-style and bridge for music. Let
$\mathcal A=\{h:w_h>0\}$. Define
\[
 \mathcal P=\begin{cases}
 \mathcal Q\cap\mathcal A,&\mathcal Q\cap\mathcal A\ne\varnothing,\\
 \mathcal Q,&\text{otherwise}.
 \end{cases}
\]
Assume $\mathcal Q$ is nonempty, as in these supplied datasets. For $h\in\mathcal P$,
let $\widetilde w_h=w_h$ if its positive weight was retained, and let
$\widetilde w_h=0.5$ otherwise. Let $z_{\rm pop}$ denote the standardized
$\log(1+f_i)$ score, with weight $w_{\rm pop}=0$ if it was not selected. Then
\begin{equation}\label{P}
 P=\sum_{h\in\mathcal P}\widetilde w_hz_h^0+\gamma z_R+w_{\rm pop}z_{\rm pop}.
\end{equation}
The channels in $\mathcal Q$ are unaffected by the incidence modification, so
$z_h^0=z_h^1$ there. The fallback retains selected popularity and Rao terms as well
as category/descriptive channels. It is neither independent of history nor disjoint
from $Y$. Its restricted channel set is a modeling choice motivated by coverage,
not a theorem requiring this exact set.

\section{Frequency-dependent backoff: an exact interpretation}
Let $s_0\geq0$ be a validation-selected fallback scale and set $\mu_0=s_0P$.
Let $\kappa\geq0$ be the selected backoff strength. Define the final score by
\begin{equation}\label{theta}
 \theta(c\mid S)=\begin{cases}
 Y(c\mid S),&\kappa=0,\\[2pt]
 \displaystyle\frac{f_cY(c\mid S)+\kappa\mu_0(c\mid S)}{f_c+\kappa},&\kappa>0.
 \end{cases}
\end{equation}
When enabled, backoff increases the fallback's relative weight for candidates with
less training history. Treating frequency as a measure of the useful amount of
history is an assumption; frequency is not a calibrated estimate of uncertainty.

\begin{proposition}[Pointwise regularization identity]
Fix $S,c,Y,\mu_0$ and $\kappa>0$. Over $t\in\mathbb R$, define
\[
 J_c(t)=f_c(t-Y)^2+\kappa(t-\mu_0)^2.
\]
The score in~\eqref{theta} is the unique minimizer of $J_c$.
\end{proposition}
\begin{proof}
$J'_c(t)=2f_c(t-Y)+2\kappa(t-\mu_0)$ and
$J''_c(t)=2(f_c+\kappa)>0$. Setting the derivative to zero gives
$t=(f_cY+\kappa\mu_0)/(f_c+\kappa)$.
\end{proof}
This identity explains how a fixed pair of scores is combined. It is not the
objective optimized by the entire fitting pipeline, and does not prove optimality
of the frequency weight, fallback construction, or ranking risk. If $f_c=\kappa=0$,
the quadratic is constant; returning $Y$ is the explicit disabled-backoff convention.

A Gaussian posterior interpretation requires additional distributional assumptions.
The archived four supplied-split fits have raw slopes approximately
$-0.273,-2.528,-5.336,-5.869$ when mean squared disagreement between $Y$ and the scaled
fallback is regressed on $1/f_c$ over warm items. Those signs do not support the
posited positive inverse-frequency variance relation in those fits. They are not
formal rejection tests, and disagreement between two scores is not ground-truth
prediction error. The procedure clips variance estimates and selects backoff settings
using validation. Thus the exact identity above is a firmer explanation than a claim
of empirically validated posterior uncertainty.

\section{Consequences of shared terms}
For $\kappa>0$, define $\rho_c=f_c/(f_c+\kappa)$ and $H=Y-P$. Substitution gives
\begin{equation}\label{sharedidentity}
 \theta=\rho_cH+\bigl[\rho_c+(1-\rho_c)s_0\bigr]P.
\end{equation}
At $s_0=1$, this reduces to $P+\rho_cH$. At a different scale, the coefficient of
$P$ also varies by candidate. $H$ is defined algebraically; it can be interpreted
as the remaining selected channels when shared coefficients agree, as in the four
archived selections. Default fallback weights require separate interpretation.

For comparison, let $a\in[0,1]$ be a candidate-independent fallback fraction. Its score is
\[
 \theta_a=(1-a)Y+as_0P=(1-a)H+(1-a+as_0)P.
\]
When $Y$ and $P$ use a common fixed signal dictionary, this constant mixture is a
reweighting of that dictionary. Improvements over a coarse, sequentially selected
weight vector do not by themselves establish a new representational capability.
Frequency-dependent backoff permits different coefficients for different candidates;
its practical advantage over constant reweighting must be measured.

\begin{proposition}[Ordering within the cold subset]
Let $\mathcal C_0=\{c\in\I\setminus S:f_c=0\}$. Hold all component scores and
weights fixed. Suppose $Y(c\mid S)-P(c\mid S)=b(S)$ for all $c\in\mathcal C_0$,
where $b(S)$ does not depend on $c$. If $s_0>0$ and $\kappa>0$, enabling backoff
preserves the ordering of candidates within $\mathcal C_0$, including ties.
\end{proposition}
\begin{proof}
For $c,d\in\mathcal C_0$, the disabled score difference is
$Y(c\mid S)-Y(d\mid S)=P(c\mid S)-P(d\mid S)$. The enabled difference is
$s_0[P(c\mid S)-P(d\mid S)]$, which has the same sign.
\end{proof}
This is a statement in exact real arithmetic. With exact scores, an identical
item-ID tie rule preserves the within-subset ranking. Floating-point evaluation can
create or break near-ties; bit-identical numerical rankings do not follow from the proof.
The proposition does not preserve rankings between cold and warm candidates, nor
apply to a comparison that also retunes components. One sufficient condition is
that $Y-P$ contains only standardized history channels whose raw scores are zero
for every cold candidate. Catalogue standardization then contributes the same
constant for each cold candidate within each such channel. Auxiliary histories,
different fallback coefficients, or $s_0=0$ can invalidate the premises.

\section{Fitting and the empirical argument}\label{evidence}
Weights are selected by a finite greedy coordinate search on validation ranking
quality, followed by optional-term and backoff selection. This is a fitted procedure
with order dependence and finite search resolution; it is not guaranteed to find
a global optimum. In the mechanism study, every reported recipe is frozen before
test scoring. Catalogue features are available transductively. No conclusion here
asserts that external feature encoders were trained without evaluation-item exposure.

The completed studies yield the following Recall@20 values:
\begin{center}
\begin{tabular}{@{}lrr@{}}\toprule
Construction & POG & POG\_dense \\\midrule
Historical family & 0.030875 & 0.375379\\
Shared category/descriptive family & 0.044542 & 0.112917\\
Combined, without refinements & 0.051458 & 0.390928\\
Full BundleFuse & 0.054458 & 0.395223\\\bottomrule
\end{tabular}
\end{center}
The historical family includes popularity. The family-only models are separately
validation-selected; they are not frozen-weight removals. The combined model exceeds
both families individually in Recall@20 and NDCG@20 under the current correction
reserving all 96 planned comparisons. This supports complementary predictive
information in the tested constructions. It does not prove every channel necessary
or attribute the contrast between datasets solely to density.

With matched fallback scale, frequency backoff exceeds the selected constant mixture
on POG: Recall@20 is $0.054458$ versus $0.051833$, with corrected evidence in both
metrics. Under unit scale, that superiority is not established. On POG\_dense,
the corresponding matched-scale recalls are $0.392893$ and $0.394353$; their difference
is not established. Absence of a significant difference is not an equivalence result.

On POG, cold-target recall in full-catalogue retrieval increases from $0.039884$
without backoff to $0.044069$ with matched-scale frequency backoff, while within-cold
Recall@20 remains $0.075222$. These are descriptive bucket diagnostics, conditional
on queries containing cold targets. They agree with the ordering proposition's
predicted distinction. The evaluated core variants on POG\_dense retrieve no cold
targets in the global top 20, despite nonzero within-cold recall; universal cold-start
success is therefore unsupported.

The five tested selection conditions give Recall@20 ranges
$[0.050958,0.054458]$ on POG and $[0.385736,0.395896]$ on POG\_dense.
These descriptive ranges show that selection sensitivity is material. The best
variant must not be chosen using test results. The studies are follow-up analyses
on previously inspected benchmarks, with uncertainty conditional on fitted models
and fixed splits. The final joint correction across datasets awaits the music runs.

\section{The resulting publication claim}
The supported central argument is that explicit historical estimates and evidence
shared through item attributes improve one another in the completed experiments.
Backoff and multi-item corrections extend that construction, with benefits that
depend on the evaluated conditions. The contribution must concern the particular
score construction and its demonstrated behavior relative to appropriate controls.
It cannot rest solely on combining information sources, avoiding a neural encoder,
or rewriting a convex mixture as a quadratic minimizer.

This positioning is consistent with prior work: EASE establishes strong explicit
item-relation modeling~\cite{ease}; CLHE and RaMen already combine several information
sources for bundle construction~\cite{clhe,ramen}; empirical-Bayes product search
already combines behavioral estimates with descriptive priors~\cite{eb}.
The complete related-work and benchmark-provenance comparison remains necessary.

\paragraph{Audit trail.}
The equations were checked against the scoring and fusion implementation at commit
\texttt{80ff927}. The new POG results come from
the dataset-level \texttt{summary.json} files under
\texttt{results/publication\_mechanisms};
the four historical variance fits come from the supplied-split selections at
\texttt{4516eab}. Historical and current experiment records are distinguished.
An independent synthetic-catalogue audit checked 52 equation cases against the
implementation, including fallback-default weights, propagation, incidence, Rao
scores, and enabled/disabled backoff. The largest absolute discrepancy was
$4.77\times10^{-7}$, consistent with the tested float32 computations. Numerical
agreement checks implementation correspondence, not statistical optimality.
The audit is retained in \texttt{results/argument\_audit/equations.json}.
The method source and running experiments were not changed by this document revision.

\begin{thebibliography}{9}\small
\bibitem{ease} H. Steck. Embarrassingly Shallow Autoencoders for Sparse Data. WWW, 2019.
\url{https://arxiv.org/abs/1905.03375}.
\bibitem{clhe} Y. Ma et al. Leveraging Multimodal Features and Item-level User Feedback
for Bundle Construction. WSDM, 2024. \url{https://arxiv.org/abs/2310.18770}.
\bibitem{ramen} H.-S. Nguyen et al. RaMen: Multi-Strategy Multi-Modal Learning for Bundle
Construction. 2025 preprint. \url{https://arxiv.org/abs/2507.14361}.
\bibitem{eb} C. Han et al. Addressing Cold Start in Product Search via Empirical Bayes.
CIKM, 2022. \url{https://doi.org/10.1145/3511808.3557066}.
\end{thebibliography}
\end{document}
